% glossary.tex -- shared by the analysis note and the technical supplement.
% Plain-language definitions, in the order a reader meets the ideas, not alphabetical.
\section*{Glossary}
\addcontentsline{toc}{section}{Glossary}
\label{sec:glossary}

\begin{description}\itemsep2pt
  \item[Signal definition] The set of particles leaving the nucleus that counts as the process being
    measured, with momentum thresholds and a fiducial requirement. It is deliberately \emph{not} a
    statement about the interaction mechanism, because final-state interactions mix mechanisms and
    only the final state is observable.
  \item[Fake data / pseudo-data] A simulated dataset used in place of beam-on data while the
    analysis is blind. Here it is a Poisson throw: each simulated event is replicated a
    Poisson-distributed number of times with mean equal to its weight at data exposure, so the
    sample has the statistical fluctuations of a real dataset.
  \item[Blind] The signal region has not been looked at in beam-on data. Every plotted ``data''
    point in these documents is fake data.
  \item[Beam-on, beam-off (EXT), dirt] Beam-on data are recorded in the beam spill; beam-off data are
    recorded between spills and measure cosmic-ray activity; ``dirt'' is simulated neutrino
    interactions outside the cryostat whose products enter the detector. The prediction is
    neutrino simulation $+$ beam-off $+$ dirt.
  \item[Gate ratio] The factor scaling beam-off data to the beam-on exposure: the ratio of the number
    of beam-on to beam-off trigger gates in a running period.
  \item[Tune, central value (CV)] The \uboone{} tune is the collaboration's adjusted GENIE model; the
    central value is the nominal prediction from it. ``CV weight'' is the per-event weight that
    turns the raw simulation into the tuned prediction.
  \item[Safe-weight rule] Any event weight that is non-finite, negative or above $30$ is replaced by
    $1$. One rule, applied everywhere, so that the prediction and every pseudo-data throw agree.
  \item[Universe] One alternative version of the prediction with a systematic parameter varied. A
    \emph{multisim} is a set of hundreds of universes drawn randomly from the parameter
    uncertainties; a \emph{unisim} is a single $\pm1\sigma$ variation of one parameter. PPFX
    universes vary the flux; GENIE universes the interaction model; reinteraction universes the
    hadron transport.
  \item[Detector variations (detVar)] Dedicated simulation samples with one aspect of the detector
    response changed --- light yield, recombination, space-charge, wire response --- used to
    estimate the detector uncertainty.
  \item[Migration matrix, efficiency, smearceptance] The migration matrix $M_{ij}$ is the probability
    that an event in true bin $j$ is reconstructed in bin $i$; the efficiency $\varepsilon_j$ is the
    fraction of true bin $j$ that is selected at all. Their product $S_{ij}=\varepsilon_j M_{ij}$
    is the response the unfolding inverts; we call it the smearceptance matrix.
  \item[Migration diagonal] The fraction of events reconstructed in a bin that were truly in that
    bin (the column-normalised diagonal of $M$). A binning is accepted only if every bin's diagonal
    exceeds $0.68$.
  \item[Unfolding, Wiener-SVD] Solving for the true spectrum given the reconstructed one and the
    response. A plain inversion amplifies statistical noise; Wiener-SVD regularises it with a filter
    built from the simulation, trading a small bias for a large reduction in variance.
  \item[Additional-smearing matrix $A_C$] The price of regularisation, made explicit: the unfolded
    result estimates $A_C\,T$ rather than the truth $T$. A prediction must be multiplied by $A_C$
    before it is compared with the result. $A_C$ is released with every extraction.
  \item[Not normalisation-preserving] The rows of $A_C$ do not sum to one, so the integral of an
    unfolded differential result is not the total cross section and differs between observables of
    the same sample. The total is measured separately.
  \item[One-bin extraction] The total cross section, extracted with a single true bin covering the
    whole phase space. Nothing is regularised, $A_C=1$, and the result is the physical total.
  \item[Conditioning, $s_{\min}/s_1$, rank-deficient] The ratio of the smallest to the largest
    singular value of the regularised response. Near zero, some direction in the true spectrum
    cannot be recovered from the data at all (rank-deficient); the result in that direction is the
    prior. A binning can pass the migration-diagonal criterion and still be unusable here.
  \item[Closure] Unfolding pseudo-data and checking that the truth comes back. Closure against the
    $A_C$-smeared truth tests the implementation; it cannot test whether the model is right,
    because the same model is on both sides.
  \item[Covariance] The matrix of uncertainties and their correlations between bins, one per
    systematic source, propagated through the unfolding. ``Nominal configured covariance'' means the
    sum over the sources implemented in the analysis; the sources known to be missing are listed
    in the limitations.
  \item[Flux-averaged cross section] The cross section averaged over the neutrino energy spectrum
    of the beam, per argon nucleus, in the stated phase space. It depends on the flux, which is why
    each horn configuration is a separate measurement.
  \item[FHC, RHC] Forward and reverse horn current: the two polarities of the NuMI focusing horns,
    giving a neutrino-dominated and an antineutrino-enriched beam.
  \item[Beam frame] All angles are measured about the neutrino direction, which at \uboone{}'s
    off-axis position is $28^\circ$ from the detector $z$ axis.
  \item[Control region, sideband] A sample selected by inverting one requirement of the signal
    selection, so that it is disjoint from the signal region and enriched in one background. It
    can be inspected on beam-on data while the signal region stays blind.
  \item[Transverse kinematic imbalance (TKI)] Variables built from the momenta transverse to the
    neutrino direction of the lepton and the hadronic system: $\delta p_T$, $\delta\alpha_T$,
    $\delta\phi_T$ and the inferred nucleon momentum $p_n$. They isolate nuclear effects from the
    unknown neutrino energy.
  \item[LLR particle identification] The log-likelihood-ratio score built from track $dE/dx$ that
    separates protons from muons and pions.
  \item[Generator predictions (FTE files)] Standalone GENIE, NuWro, NEUT and GiBUU predictions,
    run with the NuMI flux and binned identically to the results, stored as histograms for the
    comparison figures.
\end{description}
